\documentclass[11pt,a4paper]{article}

% ============================================================================
% PACKAGES
% ============================================================================
\usepackage[utf8]{inputenc}
\usepackage[T1]{fontenc}
\usepackage{amsmath,amssymb,amsthm}
\usepackage{graphicx}
\usepackage{booktabs}
\usepackage{siunitx}
\usepackage{algorithm}
\usepackage{algpseudocode}
\usepackage[hyphens]{url}
\usepackage[colorlinks=true,citecolor=blue,linkcolor=blue,urlcolor=blue]{hyperref}
\usepackage[backend=biber,style=numeric,sorting=none]{biblatex}
\usepackage{caption}
\usepackage{subcaption}
\usepackage{multirow}
\usepackage{array}
\usepackage{xcolor}

% ============================================================================
% BIBLIOGRAPHY
% ============================================================================
\addbibresource{references.bib}

% ============================================================================
% FORMATTING
% ============================================================================
\setlength{\parindent}{0pt}
\setlength{\parskip}{6pt}

\captionsetup{font=small,labelfont=bf}

% ============================================================================
% METADATA
% ============================================================================
\title{\textbf{Shape-First Behavioral Segmentation of Household Energy Consumption: A Reproducible Framework with Evidence-Based K-Selection and Controlled Validation}}

\author{
  [Author Name]$^{1}$ \and
  [Author Name]$^{1}$ \and
  [Author Name]$^{1}$ \and
  [Author Name]$^{1}$ \and
  [Author Name]$^{1}$\\[6pt]
  \small $^{1}$[Institution/Department]\\
  \small \texttt{\{email addresses\}}
}

\date{}

% ============================================================================
% DOCUMENT
% ============================================================================
\begin{document}

\maketitle

% ============================================================================
% ABSTRACT
% ============================================================================
\begin{abstract}
Household energy consumption segmentation typically groups consumers by total usage magnitude, overlooking temporal behavioral patterns critical for demand response and grid management. We present a reproducible framework that separates consumption magnitude from temporal usage behavior through shape-normalized feature engineering, enabling behavioral clustering that reflects \emph{when} energy is consumed rather than \emph{how much}. The framework employs 51 scale-invariant features derived from normalized 24-hour load profiles, applies PCA to retain 95\% variance (10 components), and selects the number of clusters K via a pre-registered multi-criterion composite rule combining silhouette coefficient, Calinski-Harabasz index, Davies-Bouldin index, and multi-seed stability analysis. Validated on a controlled synthetic dataset of 200 consumers over 365 days with known ground-truth behavioral archetypes, the framework achieves K=4 with Adjusted Rand Index (ARI) 0.813 against hidden labels never used during clustering. Multi-seed stability yields mean pairwise ARI 0.995, and temporal stability across quarterly segments achieves mean ARI 0.882, demonstrating robust behavioral groupings independent of seasonal magnitude variations. A comprehensive ablation study shows magnitude-only features achieve near-zero archetype recovery (ARI $\approx$ 0.00) despite highest silhouette scores (0.52), confirming successful separation of scale from behavior. The framework includes complete reproducibility infrastructure with cryptographic configuration hashing, deterministic execution, versioned artifacts, and documented pathways for real-world data adaptation. All code, synthetic data generation, experimental protocols, and results are publicly available. This work contributes methodological rigor and reproducibility standards for behavioral energy analytics rather than novel clustering algorithms.

\textbf{Keywords:} household energy consumption, load profile clustering, behavioral segmentation, shape normalization, PCA, K-means, evidence-based K selection, reproducible research, smart meter analytics, demand response
\end{abstract}

% ============================================================================
% MAIN CONTENT
% ============================================================================
\section{Introduction}
\label{sec:introduction}

The widespread deployment of smart meters has enabled high-resolution monitoring of household electricity consumption at unprecedented scale \cite{lazar2021investigating,kwac2014household}. This granular data enables segmentation of consumers into behaviorally distinct groups, supporting targeted demand response programs, dynamic pricing strategies, and infrastructure planning \cite{yan2021multiple,abdullah2025ml}. However, traditional clustering approaches often group households primarily by total consumption magnitude---\emph{how much} energy they use---rather than by temporal usage patterns---\emph{when} energy is consumed \cite{jin2016loadshape,rhodes2014clustering}.

Magnitude-based segmentation conflates two independent dimensions: a household's energy intensity (driven by building size, occupancy, appliances) and its behavioral routine (driven by occupancy patterns, work schedules, lifestyle preferences). Two households with identical daily load shapes but different consumption scales will be separated by magnitude-driven clustering, while households with opposite peak hours but similar totals may be grouped together. For demand response applications, behavioral timing is often more actionable than aggregate consumption \cite{haben2016analysis}.

\subsection{Motivation}

Consider two households: one peaks at midday (home office, daytime HVAC), another peaks at evening (commuter, evening cooking and entertainment). If the first consumes 30 kWh/day and the second 25 kWh/day, magnitude-driven clustering may separate them despite nearly identical timing. Conversely, if both consume 30 kWh/day, they may cluster together despite opposite demand profiles. For a utility designing time-of-use rates or load-shifting incentives, the \emph{temporal pattern} is the actionable signal \cite{afzalan2020twostage,valdes2021smart}.

Separating magnitude from timing requires explicit normalization: transforming each household's 24-hour load curve into a shape representation invariant to scalar multiplication. While shape-based clustering has been explored \cite{rahayu2021shape,jin2016loadshape}, comprehensive frameworks addressing feature engineering, dimensionality reduction, evidence-based K selection, stability validation, and reproducibility remain limited in the literature.

\subsection{Research Gap}

Existing household energy clustering studies face several methodological challenges:

\begin{enumerate}
    \item \textbf{Feature Engineering:} Many approaches use raw hourly values or magnitude-contaminated features, allowing consumption scale to dominate distance calculations \cite{kwac2014household,rhodes2014clustering}.
    
    \item \textbf{K Selection:} The number of clusters K is often chosen heuristically (e.g., ``3-5 for interpretability'') or via single metrics (elbow method, silhouette alone) without systematic evaluation \cite{hennig2020comparing}.
    
    \item \textbf{Validation:} Internal metrics (silhouette, Calinski-Harabasz) measure cluster quality but cannot confirm meaningful behavioral distinctions. External validation requires ground truth rarely available in real datasets \cite{gates2019adjusted}.
    
    \item \textbf{Stability:} Cluster robustness across random initializations and temporal windows is often unreported, leaving uncertainty about partition reliability \cite{henning2020adjusting}.
    
    \item \textbf{Reproducibility:} Many studies lack publicly available code, data, or sufficient methodological detail for independent replication \cite{peng2011reproducible}.
\end{enumerate}

\subsection{Contributions}

This paper addresses these gaps through a comprehensive, reproducible framework for shape-first behavioral segmentation. Our specific contributions are:

\begin{enumerate}
    \item \textbf{Shape-Normalized Feature Engineering:} 51 scale-invariant behavioral features derived from normalized 24-hour load profiles, explicitly separating magnitude from temporal patterns (Section~\ref{sec:features}).
    
    \item \textbf{Evidence-Based K Selection:} Pre-registered multi-criterion composite rule combining silhouette, Calinski-Harabasz, Davies-Bouldin, cluster balance constraints, and multi-seed stability, with parsimony preference when solutions are statistically indistinguishable (Section~\ref{sec:kselection}).
    
    \item \textbf{Controlled Validation:} Synthetic dataset with known ground-truth behavioral archetypes enabling quantitative assessment of recovery (ARI=0.813), while explicitly hiding ground truth during all modeling steps (Section~\ref{sec:dataset}).
    
    \item \textbf{Comprehensive Stability Analysis:} Multi-seed clustering stability (mean ARI=0.995) and temporal stability across quarterly segments (mean ARI=0.882), confirming robust behavioral groupings independent of seasonal variations (Section~\ref{sec:stability}).
    
    \item \textbf{Rigorous Ablation Study:} Systematic comparison of five feature sets (magnitude-only, shape-only, behavioral, combined) demonstrating that magnitude-only features achieve near-zero behavioral recovery (ARI $\approx$ 0.00) despite highest internal quality metrics (Section~\ref{sec:ablation}).
    
    \item \textbf{Complete Reproducibility Infrastructure:} Deterministic configuration with cryptographic hashing, pinned dependencies, versioned artifacts, documented execution protocols, and publicly available implementation (Section~\ref{sec:reproducibility}).
\end{enumerate}

\subsection{Scope and Positioning}

This work is a \textbf{methodology paper} emphasizing reproducible research practices for behavioral energy analytics. We do not claim novelty in PCA or K-means algorithms (both well-established \cite{jolliffe2002pca,lloyd1982least,arthur2007kmeans}), but rather contribute rigorous implementation, comprehensive validation, and reproducibility standards. The primary experimental validation uses synthetic data with known ground truth; a documented pathway for real-world data exists but is not executed in this study (discussed in Section~\ref{sec:limitations}).

\subsection{Paper Organization}

Section~\ref{sec:related} reviews related work on household energy clustering, PCA-based dimensionality reduction, and validation methodologies. Section~\ref{sec:methodology} details the complete analytical pipeline from data generation through cluster profiling. Section~\ref{sec:experiments} describes experimental design, configurations, and evaluation protocols. Section~\ref{sec:results} presents flagship results, ablation studies, and stability analyses. Section~\ref{sec:discussion} interprets findings and practical implications. Section~\ref{sec:limitations} acknowledges methodological constraints and dataset scope. Section~\ref{sec:conclusion} summarizes contributions and future directions.

\section{Related Work}
\label{sec:related}

\subsection{Household Energy Load Profiling}

Smart meter deployment has enabled fine-grained analysis of residential electricity consumption \cite{kwac2014household,rhodes2014clustering}. Early work focused on aggregated load profiles for utility-level forecasting \cite{haben2016analysis}, but increasing attention has shifted to disaggregated household-level patterns revealing behavioral heterogeneity \cite{lazar2021investigating,jin2016loadshape}.

Kwac et al.~\cite{kwac2014household} applied hierarchical clustering to hourly smart meter data, identifying distinct consumption patterns but noting that magnitude variations dominated the groupings. Rhodes et al.~\cite{rhodes2014clustering} used K-means on raw consumption features, achieving interpretable clusters but without systematic separation of scale from shape. Jin et al.~\cite{jin2016loadshape} at Lawrence Berkeley National Laboratory developed foundational clustering methodologies for residential smart meter data, emphasizing the importance of feature engineering and validation.

Recent work increasingly recognizes the value of shape-based clustering. Lazar et al.~\cite{lazar2021investigating} applied adaptive K-means to daily load shapes, finding outdoor temperature as the dominant driver of variability in summer-peaking utilities. Afzalan et al.~\cite{afzalan2020twostage} proposed a two-stage approach separating temporal patterns from peak demand using Dynamic Time Warping, explicitly addressing shape misalignment. Vald\'es et al.~\cite{valdes2021smart} introduced two-stage k-medoids clustering of normalized load shapes for demand response in high-renewable contexts.

\subsection{PCA and Dimensionality Reduction}

Principal Component Analysis (PCA) is widely used for dimensionality reduction in energy analytics \cite{jolliffe2002pca}. Haben et al.~\cite{haben2019fpca} applied functional PCA to electric consumption patterns for multi-horizon forecasting, demonstrating PCA's effectiveness in capturing temporal load structure. Liao et al.~\cite{liao2016pattern} combined PCA with logistic regression for pattern-based energy consumption analysis, using PCA to mitigate multicollinearity among correlated features.

Ding and He~\cite{ding2004kmeans} established theoretical connections between K-means and PCA, showing that K-means cluster indicators span the principal subspace under certain conditions. This relationship motivates PCA as a preprocessing step for K-means, reducing noise and computational cost while preserving cluster-relevant variance.

\subsection{Shape Normalization and Feature Engineering}

Several studies explicitly address magnitude-timing separation. Rahayu et al.~\cite{rahayu2021shape} proposed shape-based clustering using Dynamic Time Warping to align temporal patterns independent of magnitude. Afzalan et al.~\cite{afzalan2020twostage} extracted centroids accounting for shape misalignments before merging clusters of similar magnitude. Vald\'es et al.~\cite{valdes2021smart} normalized daily time series to unit sum, enabling clustering by temporal distribution rather than absolute consumption.

However, most approaches focus on shape alignment or normalization of raw hourly sequences, without systematic engineering of scale-invariant behavioral descriptors (e.g., peak concentration, harmonic structure, weekend behavior, load variability). Our work contributes a comprehensive 51-feature behavioral representation explicitly designed for magnitude-timing separation.

\subsection{Cluster Validation and K Selection}

Determining optimal K remains a fundamental challenge in unsupervised clustering \cite{hennig2020comparing}. Common heuristics include elbow methods (inertia curve inflection), silhouette analysis \cite{rousseeuw1987silhouettes}, and variance-ratio criteria (Calinski-Harabasz \cite{calinski1974dendrite}). Davies-Bouldin index \cite{davies1979cluster} measures cluster compactness versus separation. However, single-metric approaches may yield conflicting recommendations, and internal metrics cannot confirm semantic meaningfulness \cite{gates2019adjusted}.

Hennig and Liao~\cite{hennig2020comparing} proposed aggregating multiple calibrated validity indices, recognizing that different metrics capture complementary aspects of cluster quality. Gates and Ahn~\cite{gates2019adjusted} reviewed partition comparison indices including Adjusted Rand Index (ARI), emphasizing permutation-invariance and correction for chance agreement. Henning et al.~\cite{henning2020adjusting} proposed Modified ARI (MARI) based on multinomial rather than hypergeometric models, addressing bias in dependent clusterings.

Our framework adopts a pre-registered composite rule combining multiple metrics with cluster balance and stability constraints, selecting the simplest (smallest K) solution within tolerance when multiple candidates are statistically indistinguishable.

\subsection{Stability and Robustness Analysis}

Clustering stability across random initializations is critical for reliable findings \cite{henning2020adjusting}. Pairwise ARI between partitions from different seeds quantifies reproducibility independent of ground truth. Temporal stability---assessing whether cluster structure persists across time windows---has received less attention in energy literature despite its practical importance for longitudinal demand response.

Lazar et al.~\cite{lazar2021investigating} evaluated variability drivers across different time scales but did not systematically assess clustering stability. Afzalan et al.~\cite{afzalan2020twostage} focused on within-window cluster quality without temporal validation. Our work explicitly quantifies both multi-seed stability and temporal robustness through quarterly re-analysis with consistent protocols.

\subsection{Synthetic Data for Validation}

Controlled validation with known ground truth is rarely feasible with real smart meter data due to privacy constraints and absence of behavioral labels. Synthetic datasets enable quantitative assessment of clustering recovery \cite{meinrenken2022synthetic,makonin2020synd}.

Meinrenken et al.~\cite{meinrenken2022synthetic} developed high-resolution synthetic residential energy profiles for the United States using bottom-up modeling with extensive validation against reported energy use. Makonin et al.~\cite{makonin2020synd} released SynD, 180 days of synthetic power data for non-intrusive load monitoring. Pfeifer et al.~\cite{pfeifer2024synthetic} proposed conditional generation of household load profiles for smart grid applications.

Our synthetic generator assigns each consumer to a behavioral archetype (daytime-peaking, evening-peaking, flat, weekend-heavy) while drawing magnitude independently, enabling separation validation. Critically, archetype labels are hidden during all modeling steps and used only for post-hoc recovery assessment.

\subsection{Reproducibility in Computational Research}

Reproducibility is increasingly recognized as essential for scientific rigor \cite{peng2011reproducible,stodden2013setting}. Peng~\cite{peng2011reproducible} argues that computational research must provide code, data, and execution protocols to enable independent verification. Stodden et al.~\cite{stodden2013setting} advocate for reproducibility as the default standard in computational and experimental mathematics.

In energy analytics, reproducibility remains limited. Many studies lack public code or sufficient detail for replication. Our framework contributes deterministic configuration with cryptographic hashing, pinned dependencies, versioned artifacts, and complete source code availability, aiming to set reproducibility standards for behavioral energy research.

\subsection{Positioning of This Work}

This paper synthesizes and extends prior work through:
\begin{itemize}
    \item Comprehensive 51-feature behavioral representation (beyond simple normalization \cite{valdes2021smart} or raw sequences \cite{kwac2014household})
    \item Evidence-based K selection integrating multiple metrics, stability, and parsimony \cite{hennig2020comparing}
    \item Controlled validation with synthetic ground truth hidden during modeling \cite{meinrenken2022synthetic}
    \item Systematic stability analysis across both random seeds and temporal windows
    \item Rigorous ablation demonstrating magnitude-timing separation
    \item Complete reproducibility infrastructure \cite{peng2011reproducible}
\end{itemize}

We emphasize that our contribution is methodological integration and reproducibility, not algorithmic novelty in PCA \cite{jolliffe2002pca} or K-means \cite{lloyd1982least,arthur2007kmeans}, both well-established techniques.

\section{Methodology}
\label{sec:methodology}

This section details the complete analytical pipeline from data generation through cluster profiling. Figure~\ref{fig:pipeline} illustrates the end-to-end workflow.

\subsection{Dataset Generation}
\label{sec:dataset}

\subsubsection{Synthetic Data Rationale}

Real household smart meter data with known behavioral ground truth is unavailable due to privacy constraints and absence of labels. Synthetic data enables controlled validation: each consumer is assigned a behavioral archetype, allowing quantitative assessment of clustering recovery while maintaining the independence property critical for evaluating magnitude-timing separation.

\subsubsection{Archetype-Based Generator}

The synthetic generator produces hourly electricity consumption for $N$ consumers over $D$ days. Each consumer $i$ is assigned to one of four behavioral archetypes $\mathcal{A} = \{\text{daytime}, \text{evening}, \text{flat}, \text{weekend}\}$:

\begin{itemize}
    \item \textbf{Daytime:} Peaks 12:00-14:00, lower evening demand (home office, midday HVAC)
    \item \textbf{Evening:} Peaks 18:00-21:00, lower daytime demand (commuter, evening cooking/entertainment)
    \item \textbf{Flat:} Approximately uniform 24-hour distribution (commercial-like, always-on baseline)
    \item \textbf{Weekend:} Evening-peaking with elevated weekend consumption (lifestyle-driven weekly variation)
\end{itemize}

For consumer $i$ at hour $h$ on day $d$:

\begin{equation}
E_{i,h,d} = B_i \cdot S_{\mathcal{A}_i}(h) \cdot M_i(d) \cdot (1 + \epsilon_{i,h,d})
\end{equation}

where:
\begin{itemize}
    \item $B_i \sim \text{Uniform}(0.5, 3.0)$ is base consumption magnitude (kW), drawn \emph{independently} of archetype $\mathcal{A}_i$
    \item $S_{\mathcal{A}_i}(h)$ is the normalized 24-hour shape template for archetype $\mathcal{A}_i$ (sums to 1)
    \item $M_i(d)$ is the seasonal magnitude multiplier for consumer $i$ on day $d$ (mean-corrected, independent of $\mathcal{A}_i$)
    \item $\epsilon_{i,h,d} \sim \mathcal{N}(0, \sigma^2)$ is Gaussian noise
\end{itemize}

\textbf{Critical Property:} Magnitude $B_i$ and seasonal phase are drawn independently of archetype $\mathcal{A}_i$. This construction ensures that consumption scale carries no information about behavioral pattern, enabling validation of magnitude-timing separation.

\subsubsection{Ground Truth Handling}

Archetype labels $\mathcal{A}_i$ and seasonal phases are stored in the generated dataset but \textbf{dropped before any preprocessing or modeling}. These labels are used exclusively for post-hoc validation via Adjusted Rand Index (ARI) and never influence feature engineering, PCA, or K-means. This independence is verified through dataset integrity checks (Section~\ref{sec:experiments}).

\subsection{Preprocessing}
\label{sec:preprocessing}

The preprocessing pipeline ensures data quality without information leakage:

\begin{enumerate}
    \item \textbf{Timestamp Parsing:} Robust parsing with failed-parse logging
    \item \textbf{Deduplication:} Remove exact $(i, t)$ duplicates
    \item \textbf{Sorting:} Sort by $(i, t)$ for contiguous consumer-level access
    \item \textbf{Gap Filling:} Within-consumer forward/backward fill (\emph{never} cross-consumer)
    \item \textbf{Outlier Handling:} Configurable winsorization (preserves behavioral peaks rather than deletion)
    \item \textbf{Validation:} Schema, continuity, value-range checks
\end{enumerate}

\textbf{Leakage Prevention:} Missing value imputation operates within consumer groups ($i$) only, preventing contamination across households. Ground truth columns $\mathcal{A}_i$ are explicitly dropped before any statistical computation.

\subsection{Feature Engineering}
\label{sec:features}

Feature engineering is the core contribution, separating magnitude from temporal behavior through explicit normalization.

\subsubsection{Shape Normalization}

For each consumer $i$, compute mean hourly consumption $\mu_{i,h} = \frac{1}{D} \sum_{d=1}^{D} E_{i,h,d}$ across all days. The normalized 24-hour shape is:

\begin{equation}
\tilde{S}_{i,h} = \frac{\mu_{i,h}}{\sum_{h'=0}^{23} \mu_{i,h'}}
\end{equation}

This transformation ensures $\sum_{h=0}^{23} \tilde{S}_{i,h} = 1$ and makes the shape invariant to scalar multiplication: $\alpha E_{i,h,d} \rightarrow \tilde{S}_{i,h}$ (same shape).

\subsubsection{Feature Groups}

We engineer 51 behavioral features organized into five groups:

\begin{table}[h]
\centering
\small
\caption{Behavioral Feature Groups}
\label{tab:feature_groups}
\begin{tabular}{@{}lcp{7cm}@{}}
\toprule
Group & Count & Description \\
\midrule
Shape & 24 & Normalized hourly bins $\tilde{S}_{i,0}, \ldots, \tilde{S}_{i,23}$ \\
Timing & 11 & Period shares (morning, afternoon, evening, night), peak timing (sin/cos), weekend ratio, night-day ratio \\
Shape Descriptors & 9 & Entropy, Gini, base-load share, Fourier harmonics (1-3), Haar wavelets (L1-L3) \\
Variability & 4 & Peak-to-average, coefficient of variation, skewness, kurtosis \\
Dispersion & 3 & Daily total CV, P90/median ratio, weekend CV ratio \\
\midrule
\textbf{Behavioral Total} & \textbf{51} & \textbf{All groups combined (shipped feature set)} \\
\bottomrule
\end{tabular}
\end{table}

\textbf{Period Shares:} Define four 6-hour blocks (night 0-6, morning 6-12, afternoon 12-18, evening 18-24). The evening share is:
\begin{equation}
\text{evening\_share}_i = \sum_{h=18}^{23} \tilde{S}_{i,h}
\end{equation}

\textbf{Peak Timing:} Encode peak hour $h^*_i = \arg\max_h \tilde{S}_{i,h}$ as circular coordinates to handle midnight wrap:
\begin{equation}
\text{peak\_hour\_sin}_i = \sin\left(\frac{2\pi h^*_i}{24}\right), \quad
\text{peak\_hour\_cos}_i = \cos\left(\frac{2\pi h^*_i}{24}\right)
\end{equation}

\textbf{Weekend Ratio:} Energy-based (not temporal):
\begin{equation}
\text{weekend\_ratio}_i = \frac{\text{mean}(E_{i,h,d} : \text{is\_weekend}(d))}{\text{mean}(E_{i,h,d} : \text{is\_weekday}(d))}
\end{equation}

\textbf{Shape Entropy:} Shannon entropy of normalized shape relative to flat profile:
\begin{equation}
H_i = -\sum_{h=0}^{23} \tilde{S}_{i,h} \log \tilde{S}_{i,h}, \quad
\text{shape\_entropy}_i = H_i / \log(24)
\end{equation}

\textbf{Harmonic Amplitudes:} Fourier components capture diurnal and semi-diurnal periodicities:
\begin{equation}
\text{harmonic\_k\_amplitude}_i = \left| \sum_{h=0}^{23} \tilde{S}_{i,h} \exp\left(\frac{2\pi i k h}{24}\right) \right|
\end{equation}
for $k = 1, 2, 3$.

\subsubsection{Scale Invariance Property}

All behavioral features satisfy the scale invariance property: if consumer consumption is multiplied by constant $\alpha > 0$, features remain unchanged. This is verified in automated tests (Section~\ref{sec:experiments}).

\subsection{Standardization and PCA}
\label{sec:pca}

\subsubsection{Standardization}

The 51-feature matrix $\mathbf{X} \in \mathbb{R}^{N \times 51}$ is standardized to zero mean and unit variance using \texttt{StandardScaler} \cite{pedregosa2011scikit}:
\begin{equation}
\mathbf{Z} = (\mathbf{X} - \boldsymbol{\mu}) \oslash \boldsymbol{\sigma}
\end{equation}
where $\oslash$ denotes element-wise division. Standardization prevents variance-dominant features from overwhelming PCA.

\subsubsection{PCA Application}

Apply PCA to $\mathbf{Z}$:
\begin{equation}
\mathbf{Z} = \mathbf{U} \boldsymbol{\Sigma} \mathbf{V}^\top
\end{equation}

Retain components until cumulative explained variance reaches threshold $\tau = 0.95$:
\begin{equation}
m = \min \left\{ k : \frac{\sum_{j=1}^{k} \lambda_j}{\sum_{j=1}^{51} \lambda_j} \geq \tau \right\}
\end{equation}
where $\lambda_j$ are eigenvalues (variances along principal components).

\textbf{Loadings vs. Weights:} PCA \texttt{components\_} are unit eigenvectors (weights). \emph{Loadings} are correlations between original features and component scores:
\begin{equation}
\text{loading}_{f,c} = \text{weight}_{f,c} \times \sqrt{\lambda_c}
\end{equation}

Loadings (bounded $[-1, 1]$) are used for interpretation; weights are used for projection.

\subsection{K-Means Clustering}
\label{sec:kmeans}

Apply K-means \cite{lloyd1982least} with k-means++ initialization \cite{arthur2007kmeans} on PCA scores $\mathbf{T} \in \mathbb{R}^{N \times m}$:
\begin{equation}
\min_{\{\mathbf{c}_k\}, \{C_k\}} \sum_{k=1}^{K} \sum_{i \in C_k} \| \mathbf{t}_i - \mathbf{c}_k \|^2
\end{equation}
where $C_k$ is cluster $k$, $\mathbf{c}_k$ is centroid, $\mathbf{t}_i$ is consumer $i$'s PCA score.

Configuration: $n\_init=10$ restarts, $max\_iter=300$, $tol=10^{-4}$, $random\_state=42$ for determinism.

\subsection{Evidence-Based K Selection}
\label{sec:kselection}

Rather than heuristic K choice, we apply a pre-registered multi-criterion rule evaluated for $K \in \{2, \ldots, 10\}$.

\subsubsection{Metrics Computed}

For each candidate $K$:
\begin{itemize}
    \item \textbf{Silhouette} \cite{rousseeuw1987silhouettes}: $s_K \in [-1, 1]$, higher = better separation
    \item \textbf{Calinski-Harabasz} \cite{calinski1974dendrite}: $CH_K > 0$, higher = better variance ratio
    \item \textbf{Davies-Bouldin} \cite{davies1979cluster}: $DB_K > 0$, lower = better compactness
    \item \textbf{Stability:} Refit K-means with 10 random seeds, compute mean pairwise ARI across partitions
    \item \textbf{Balance:} Minimum cluster size as fraction of $N$
\end{itemize}

\subsubsection{Selection Rule}

\begin{algorithm}[h]
\caption{Evidence-Based K Selection}
\label{alg:kselection}
\begin{algorithmic}[1]
\State \textbf{Input:} Candidate set $\mathcal{K} = \{2, \ldots, 10\}$, metrics $s_K, CH_K, DB_K, \text{ARI}_K$, balance $b_K$
\State \textbf{Parameters:} $\alpha_{balance} = 0.05$, $\alpha_{stability} = 0.60$, $\tau_{tolerance} = 0.05$

\State Filter: $\mathcal{K}_1 \leftarrow \{K \in \mathcal{K} : b_K \geq \alpha_{balance}\}$ \Comment{Reject tiny clusters}
\State Filter: $\mathcal{K}_2 \leftarrow \{K \in \mathcal{K}_1 : \text{ARI}_K \geq \alpha_{stability}\}$ \Comment{Reject unstable solutions}

\State Normalize: $\tilde{s}_K \leftarrow (s_K - \min s) / (\max s - \min s)$ over $\mathcal{K}_2$
\State Normalize: $\tilde{CH}_K \leftarrow (CH_K - \min CH) / (\max CH - \min CH)$
\State Normalize: $\tilde{DB}_K \leftarrow (DB_{\max} - DB_K) / (DB_{\max} - DB_{\min})$ \Comment{Invert: lower is better}

\State Composite: $\mathcal{C}_K \leftarrow \frac{1}{3}(\tilde{s}_K + \tilde{CH}_K + \tilde{DB}_K)$ for $K \in \mathcal{K}_2$

\State $K^* \leftarrow \arg\max_{K} \mathcal{C}_K$
\State $\mathcal{K}_{tol} \leftarrow \{K \in \mathcal{K}_2 : \mathcal{C}_K \geq \mathcal{C}_{K^*} - \tau_{tolerance}\}$ \Comment{Within tolerance}

\State \textbf{Return:} $\min \mathcal{K}_{tol}$ \Comment{Simplest among tied solutions}
\end{algorithmic}
\end{algorithm}

\textbf{Rationale:} Filters ensure non-degenerate, reproducible partitions. Composite score integrates complementary quality signals. Tolerance band with parsimony preference selects simplest statistically-indistinguishable solution.

\subsection{Validation Metrics}
\label{sec:validation}

\subsubsection{Internal Validation (No Ground Truth)}

\begin{itemize}
    \item \textbf{Silhouette} \cite{rousseeuw1987silhouettes}: $s_i = (b_i - a_i) / \max(a_i, b_i)$ where $a_i$ is mean intra-cluster distance, $b_i$ is mean nearest-cluster distance
    \item \textbf{Calinski-Harabasz} \cite{calinski1974dendrite}: Between-cluster to within-cluster variance ratio
    \item \textbf{Davies-Bouldin} \cite{davies1979cluster}: Average similarity of each cluster to its most similar cluster
\end{itemize}

\subsubsection{External Validation (Synthetic Data Only)}

\begin{itemize}
    \item \textbf{Adjusted Rand Index (ARI)} \cite{gates2019adjusted}: Corrects Rand Index for chance agreement, ranges $[-1, 1]$, perfect agreement = 1
    \item \textbf{Normalized Mutual Information (NMI)}: Information-theoretic similarity, ranges $[0, 1]$
\end{itemize}

\textbf{Critical:} ARI and NMI are computed \emph{after} clustering against hidden archetype labels, never used during modeling.

\subsection{Stability Analysis}
\label{sec:stability}

\subsubsection{Multi-Seed Stability}

For selected $K$, refit K-means with $R=10$ random seeds $\{r_1, \ldots, r_R\}$. Compute pairwise ARI between all partition pairs:
\begin{equation}
\text{Stability}_K = \frac{2}{R(R-1)} \sum_{i < j} \text{ARI}(\mathcal{P}_{r_i}, \mathcal{P}_{r_j})
\end{equation}

High mean ARI ($> 0.95$) indicates robust partition independent of initialization.

\subsubsection{Temporal Stability (Longitudinal Analysis)}

Partition observation window into $Q$ non-overlapping segments. Within each segment $q$:
\begin{enumerate}
    \item Re-compute all 51 behavioral features
    \item Re-apply standardization and PCA (95\% threshold)
    \item Re-run K-selection algorithm (may yield different $K_q$)
    \item Compare segment partition with full-window partition via ARI
\end{enumerate}

\begin{equation}
\text{Temporal\_Stability} = \frac{1}{Q} \sum_{q=1}^{Q} \text{ARI}(\mathcal{P}_q, \mathcal{P}_{full})
\end{equation}

High temporal ARI indicates behavioral groupings persist across seasons despite magnitude variations.

\subsection{Explainability}
\label{sec:explainability}

K-means is unsupervised and has no native feature importance. We employ post-hoc explainability:

\begin{enumerate}
    \item Train random forest classifier $f: \mathbf{X} \rightarrow \{1, \ldots, K\}$ predicting recovered cluster labels from behavioral features
    \item Evaluate surrogate via cross-validation (balanced accuracy)
    \item Apply SHAP TreeExplainer \cite{lundberg2017shap} to surrogate
\end{enumerate}

\textbf{Interpretation:} SHAP values indicate which features most distinguish clusters in the surrogate. Cross-validation accuracy is the \emph{ceiling} on explanation quality---cannot explain more than the surrogate captures.

\textbf{Critical:} Surrogate never feeds back into clustering; it is purely diagnostic.

\subsection{Cluster Profiling}

For each cluster $k$, compute statistics in \emph{original feature space}:
\begin{itemize}
    \item Mean hourly shape: $\bar{S}_{k,h} = \frac{1}{|C_k|} \sum_{i \in C_k} \tilde{S}_{i,h}$
    \item Period shares: $\sum_{h \in \text{period}} \bar{S}_{k,h}$
    \item Peak hour: $\arg\max_h \bar{S}_{k,h}$
    \item Weekend ratio, peak-to-average, CV (mean across cluster members)
\end{itemize}

Profiles are compared against population baselines to identify deviations. Human-interpretable cluster names are assigned based on dominant characteristics (e.g., ``Evening-Peaking Weekend-Heavy'').

\section{Experimental Design}
\label{sec:experiments}

\subsection{Dataset Configuration}

The flagship experimental configuration (hash \texttt{99c7a6631340d301}) comprises:
\begin{itemize}
    \item Consumers: $N = 200$
    \item Observation window: 365 days (2024-01-01 to 2024-12-30)
    \item Sampling: Hourly measurements
    \item Total records: 1,752,000 (after preprocessing)
    \item Hidden archetypes: 4 (daytime, evening, flat, weekend)
    \item Random seed: 42 (deterministic)
\end{itemize}

A reference 30-day configuration (hash \texttt{6896387297178841}) is maintained for short-horizon comparison.

\subsection{Experimental Protocol}

\subsubsection{Single-Run Flagship Experiment}

The primary evaluation uses the 365-day configuration with complete pipeline execution:
\begin{enumerate}
    \item Generate synthetic panel with hidden archetypes
    \item Validate dataset integrity (archetype-magnitude independence)
    \item Drop ground truth columns
    \item Preprocess (within-consumer imputation, no leakage)
    \item Engineer 51 behavioral features
    \item Standardize and apply PCA (95\% threshold)
    \item Sweep K=2 to 10 with multi-metric evaluation
    \item Select K via Algorithm~\ref{alg:kselection}
    \item Profile clusters in original feature space
    \item Validate against hidden archetypes (ARI, NMI)
    \item Assess multi-seed stability (10 restarts)
    \item Assess temporal stability (4 quarterly segments)
    \item Apply SHAP explainability
\end{enumerate}

\subsubsection{Ablation Study}

To validate magnitude-timing separation, we compare five feature sets on a 30-day window:
\begin{itemize}
    \item \textbf{scale}: 7 magnitude features (mean, max, sum energy, current)
    \item \textbf{shape}: 24 normalized hourly bins only
    \item \textbf{summary}: 27 derived behavioral scalars
    \item \textbf{behavioral}: 51 features (shape + summary, \textbf{shipped})
    \item \textbf{combined}: 58 features (behavioral + scale)
\end{itemize}

All arms use identical data, seed, and K-selection rule. Only feature set varies.

\subsubsection{Seed Robustness Study}

To assess generalizability, we repeat the ablation across 20 independent datasets (seeds 1-19, 42), each 30-day window with 200 consumers. This evaluates:
\begin{itemize}
    \item Mean and variance of archetype recovery per feature set
    \item Frequency of K-selection choices
    \item Statistical significance of feature set differences
\end{itemize}

Paired permutation tests assess whether feature sets differ significantly.

\subsection{Implementation}

\subsubsection{Software Stack}

\begin{itemize}
    \item Python 3.11
    \item pandas 3.0.0, numpy 2.3.5
    \item scikit-learn 1.9.0, scipy 1.18.0
    \item matplotlib 3.10.8, seaborn 0.13.2
    \item SHAP (TreeExplainer for post-hoc explanation)
\end{itemize}

All dependencies pinned in \texttt{requirements.txt} for reproducibility.

\subsubsection{Reproducibility Mechanisms}

\begin{itemize}
    \item Configuration hashing: Cryptographic hash of all parameters
    \item Deterministic execution: Fixed random seeds throughout
    \item Artifact versioning: Timestamped outputs with metadata
    \item Contract versioning: JSON artifacts with semantic versioning (v1.0.0)
\end{itemize}

\subsection{Evaluation Metrics}

See Section~\ref{sec:validation} for metric definitions. Summary:
\begin{itemize}
    \item \textbf{Internal:} Silhouette, Calinski-Harabasz, Davies-Bouldin, stability ARI
    \item \textbf{External (synthetic only):} ARI, NMI vs. hidden archetypes
    \item \textbf{Temporal:} Segment-wise ARI vs. full window
\end{itemize}

\subsection{Computational Environment}

Experiments executed on standard desktop hardware. Flagship run completion time: approximately 15 minutes for full pipeline including all stability analyses. Ablation study (5 arms × 20 seeds): approximately 2 hours total.

\subsection{Dataset Integrity Validation}

Before modeling, we verify:
\begin{enumerate}
    \item Mean consumer shapes match archetype templates
    \item Weekend effects present in weekend-archetype consumers
    \item Magnitude $B_i$ uncorrelated with archetype label (Pearson $r < 0.05$)
\end{enumerate}

This confirms the synthetic generator produces the intended structure with independent magnitude-behavior channels.

\section{Results}
\label{sec:results}

This section presents results from the flagship 365-day experiment (configuration hash \texttt{99c7a6631340d301}), ablation study, and robustness analyses. All numerical values are traceable to \texttt{outputs/reports/analysis\_summary.md} and \texttt{models/analysis\_metadata.json}.

\subsection{PCA Dimensionality Reduction}

Applying PCA to the standardized 51-feature matrix with 95\% cumulative variance threshold yields:

\begin{table}[h]
\centering
\caption{PCA Results (Flagship Configuration)}
\label{tab:pca_results}
\begin{tabular}{@{}lcc@{}}
\toprule
Parameter & Value & Notes \\
\midrule
Input features & 51 & Behavioral features \\
Components retained & 10 & Threshold: 95\% variance \\
Cumulative variance & 95.05\% & Target: 95\% \\
Kaiser criterion & 7 components & Eigenvalue $> 1$ \\
Scree elbow & 7 components & Visual inspection \\
\bottomrule
\end{tabular}
\end{table}

The 10-component solution balances dimensionality reduction with information preservation. Figure~\ref{fig:pca_variance} shows explained and cumulative variance curves.

\textbf{Top PC Interpretations (via loadings):}
\begin{itemize}
    \item \textbf{PC1} (profile shape): \texttt{profile\_ramp} +0.92, \texttt{peak\_concentration} +0.88, \texttt{harmonic\_2\_amplitude} +0.88
    \item \textbf{PC2} (day/night balance): \texttt{night\_day\_ratio} +0.92, \texttt{afternoon\_share} $-0.87$
    \item \textbf{PC3} (morning pattern): \texttt{hour\_7\_shape} +0.87, \texttt{hour\_8\_shape} +0.87
    \item \textbf{PC4} (evening timing): \texttt{hour\_17\_shape} +0.56, \texttt{hour\_18\_shape} +0.54
    \item \textbf{PC5} (weekend behavior): \texttt{weekend\_ratio} $-0.51$, \texttt{haar\_detail\_l3} +0.54
\end{itemize}

These interpretations confirm PCA captures complementary behavioral dimensions: overall profile shape, diurnal balance, specific hour patterns, and weekend effects.

\subsection{K Selection}

Algorithm~\ref{alg:kselection} applied to candidates K=2 to 10 yields complete sweep results (Table~\ref{tab:k_sweep}).

\begin{table}[h]
\centering
\small
\caption{K-Selection Sweep (Flagship, 365 days, 200 consumers)}
\label{tab:k_sweep}
\begin{tabular}{@{}cccccc@{}}
\toprule
K & Inertia & Silhouette & CH & DB & Stability ARI \\
\midrule
2 & 7083.5 & 0.2939 & 73.0 & 1.3823 & 0.9958 \\
3 & 4968.5 & 0.3305 & 93.7 & 1.1957 & 0.9852 \\
\textbf{4} & \textbf{3911.1} & \textbf{0.3283} & \textbf{96.6} & \textbf{1.1691} & \textbf{0.9947} \\
5 & 3466.1 & \textbf{0.3352} & 87.6 & 1.2023 & 0.9587 \\
6 & 3072.5 & 0.3238 & 83.6 & 1.2326 & 0.9909 \\
7 & 2844.6 & 0.3164 & 77.5 & 1.2094 & 0.8931 \\
8 & 2670.3 & 0.3072 & 72.2 & 1.2950 & 0.8651 \\
9 & 2490.6 & 0.3111 & 69.1 & 1.2015 & 0.8156 \\
10 & 2361.6 & 0.2760 & 65.6 & 1.3180 & 0.7581 \\
\bottomrule
\end{tabular}
\end{table}

\textbf{Selection Outcome:} K=4 selected with composite score 0.9444. K=5 has marginally higher silhouette (0.3352 vs 0.3283) but lower composite (0.8210). No other K within 0.05 tolerance, so parsimony rule not invoked. K=7-10 rejected for producing sub-5\% clusters.

\textbf{Inertia Elbow:} Visual elbow also at K=4, providing independent confirmation.

\subsection{Cluster Characteristics}

Table~\ref{tab:cluster_profiles} summarizes the four recovered clusters in original feature space.

\begin{table}[h]
\centering
\small
\caption{Cluster Profiles (K=4, Flagship)}
\label{tab:cluster_profiles}
\begin{tabular}{@{}lcccc@{}}
\toprule
Characteristic & Cluster 0 & Cluster 1 & Cluster 2 & Cluster 3 & Population \\
\midrule
Size (N) & 39 & 52 & 47 & 62 & 200 \\
Share (\%) & 19.5 & 26.0 & 23.5 & 31.0 & 100 \\
Peak hour & 13:00 & 19:00 & 20:00 & 19:00 & --- \\
Evening share & 0.212 & 0.258 & \textbf{0.380} & 0.299 & 0.290 \\
Afternoon share & \textbf{0.357} & 0.271 & 0.229 & 0.310 & 0.290 \\
Weekend ratio & \textbf{0.747} & 0.952 & 1.026 & \textbf{1.313} & 1.041 \\
Peak-to-avg & 8.83 & \textbf{4.92} & \textbf{11.32} & 9.45 & 8.59 \\
CV & 0.620 & \textbf{0.302} & 0.705 & 0.596 & 0.550 \\
\midrule
Assigned name & Midday- & Flat & Evening- & Evening- & --- \\
 & Peaking & All-Day & Peaking & Peaking & \\
 & Weekday- &  &  & Weekend- & \\
 & Heavy &  &  & Heavy & \\
\bottomrule
\end{tabular}
\end{table}

\textbf{Observations:}
\begin{itemize}
    \item \textbf{Cluster 0} (Midday-Peaking Weekday-Heavy): Afternoon peak (13:00), elevated afternoon share (0.357), depressed weekend ratio (0.747)
    \item \textbf{Cluster 1} (Flat All-Day): Lowest peak-to-average (4.92), lowest CV (0.302), near-uniform distribution
    \item \textbf{Cluster 2} (Evening-Peaking): Highest evening share (0.380), highest peak-to-average (11.32), concentrated evening demand
    \item \textbf{Cluster 3} (Evening-Peaking Weekend-Heavy): Evening peak with elevated weekend ratio (1.313)
\end{itemize}

\textbf{Mean Consumption Context:} Mean kWh/record for clusters [1.30, 1.38, 1.38, 1.32] vs population 1.35---approximately uniform, confirming magnitude does not drive clustering.

\subsection{Ground Truth Recovery (Synthetic Validation)}

Comparing recovered clusters against hidden behavioral archetypes (never used during modeling):

\begin{table}[h]
\centering
\caption{External Validation: Recovery by K}
\label{tab:recovery_by_k}
\begin{tabular}{@{}cccc@{}}
\toprule
K & ARI & NMI & Silhouette \\
\midrule
2 & 0.288 & 0.457 & 0.294 \\
3 & 0.602 & 0.680 & 0.331 \\
\textbf{4} & \textbf{0.813} & \textbf{0.828} & 0.328 \\
5 & 0.765 & 0.802 & \textbf{0.335} \\
6 & 0.753 & 0.782 & 0.324 \\
7-10 & $<0.74$ & $<0.78$ & $<0.32$ \\
\bottomrule
\end{tabular}
\end{table}

\textbf{Key Finding:} ARI peaks at K=4 (0.813), exactly where evidence-based rule selected. This independent validation confirms the selection was correct.

\textbf{Confusion Matrix at K=4:}

\begin{table}[h]
\centering
\caption{Archetype $\times$ Cluster Crosstab}
\label{tab:crosstab}
\begin{tabular}{@{}lcccc|c@{}}
\toprule
Archetype & Cluster 0 & Cluster 1 & Cluster 2 & Cluster 3 & Total \\
\midrule
daytime & \textbf{39} & 1 & 0 & 10 & 50 \\
evening & 0 & 0 & \textbf{47} & 3 & 50 \\
flat & 0 & \textbf{50} & 0 & 0 & 50 \\
weekend & 0 & 1 & 0 & \textbf{49} & 50 \\
\midrule
Total & 39 & 52 & 47 & 62 & 200 \\
\bottomrule
\end{tabular}
\end{table}

Three archetypes recover near-perfectly (evening 47/50, flat 50/50, weekend 49/50). Daytime archetype splits: 39/50 in Cluster 0, 10/50 in Cluster 3, reflecting behavioral overlap between midday-peaking and evening-peaking patterns when weekend behavior differs.

\subsection{Multi-Seed Stability}

Refitting K-means at K=4 with 10 random seeds yields:
\begin{itemize}
    \item Mean pairwise ARI: \textbf{0.9947 $\pm$ 0.0071}
    \item Min ARI: 0.9736
    \item Assignment agreement: 0.998
\end{itemize}

\textbf{Interpretation:} Partition is highly robust across initializations. Near-perfect stability (ARI $>$ 0.99) indicates clustering captures real structure, not random artifacts.

\subsection{Temporal Stability (Longitudinal Analysis)}

Quarterly segments re-analyzed independently:

\begin{table}[h]
\centering
\caption{Temporal Stability Across Quarters}
\label{tab:temporal_stability}
\begin{tabular}{@{}lccc@{}}
\toprule
Segment & Period & Selected K & ARI vs Full Window \\
\midrule
Q1 & 2024-01-01 to 2024-04-01 & 4 & 0.838 \\
Q2 & 2024-04-01 to 2024-07-01 & 4 & 0.892 \\
Q3 & 2024-07-01 to 2024-09-30 & 4 & \textbf{0.946} \\
Q4 & 2024-09-30 to 2024-12-30 & 4 & 0.851 \\
\midrule
\multicolumn{3}{l}{\textbf{Mean Temporal Stability}} & \textbf{0.882} \\
\bottomrule
\end{tabular}
\end{table}

\textbf{Observations:}
\begin{itemize}
    \item All segments independently select K=4, confirming evidence-based rule consistency
    \item High temporal ARI ($>0.84$ all quarters, mean 0.882) demonstrates consumer groups persist across seasons
    \item Q3 (summer) shows strongest alignment (0.946); Q1 (winter) weakest (0.838) but still strong
\end{itemize}

\textbf{Interpretation:} Behavioral groupings are stable properties of consumers, not seasonal artifacts. Despite magnitude variations (monthly mean kWh: Jan 26.1 $\rightarrow$ Jun 39.4 $\rightarrow$ Dec 25.3), temporal patterns remain coherent.

\subsection{Seasonal Analysis}

Seasonal model separately estimates magnitude and timing channels:

\begin{itemize}
    \item \textbf{Magnitude amplitude} (fractional daily swing): 0.202
    \item \textbf{Phase correlation} (vs hidden seasonal phase): Pearson $r$ = 0.678
    \item \textbf{Peak-season agreement}: 88.5\%
\end{itemize}

\textbf{Mean daily kWh by season:} Winter 26.6, Spring 35.2, Summer 38.0, Autumn 29.4

\textbf{Peak hour by season:} Autumn/Winter 19:00, Spring/Summer 20:00

\textbf{Interpretation:} Seasonal effects primarily affect consumption magnitude (amplitude 0.20) rather than timing. Phase estimation achieves moderate correlation (0.68) with hidden truth, but phase was drawn independently of archetype, so timing shifts do not create spurious clusters.

\subsection{Explainability (SHAP Analysis)}

Post-hoc random forest surrogate achieves cross-validated balanced accuracy 0.985, indicating features strongly distinguish clusters. SHAP attribution (top 3 features per cluster):

\begin{table}[h]
\centering
\small
\caption{Cluster-Discriminating Features (SHAP)}
\label{tab:shap_features}
\begin{tabular}{@{}lp{8cm}@{}}
\toprule
Cluster & Top 3 Features (SHAP importance) \\
\midrule
0 (Midday Weekday) & \texttt{afternoon\_share} (0.142), \texttt{hour\_14\_shape} (0.109), \texttt{hour\_2\_shape} (0.076) \\
1 (Flat) & \texttt{base\_load\_share} (0.175), \texttt{coefficient\_of\_variation} (0.091), \texttt{shape\_entropy} (0.060) \\
2 (Evening) & \texttt{evening\_share} (0.126), \texttt{peak\_concentration} (0.123), \texttt{harmonic\_2\_amplitude} (0.107) \\
3 (Evening Weekend) & \texttt{weekend\_ratio} (0.250), \texttt{peak\_hour\_cos} (0.091), \texttt{peak\_hour\_sin} (0.049) \\
\bottomrule
\end{tabular}
\end{table}

\textbf{Interpretation:} Each cluster distinguished by semantically meaningful features. Flat cluster driven by uniformity descriptors (base load, low CV, entropy). Weekend cluster strongly distinguished by weekend ratio (0.250 importance), confirming behavioral naming is grounded.

\subsection{Ablation Study}

Five feature sets compared on 30-day window with identical protocol:

\begin{table}[h]
\centering
\small
\caption{Ablation Study: Feature Set Comparison}
\label{tab:ablation}
\begin{tabular}{@{}lccccc@{}}
\toprule
Feature Set & Features & K & Silhouette & Stability ARI & Archetype ARI \\
\midrule
scale & 7 & 2 & \textbf{0.521} & 0.983 & \textbf{-0.004} \\
shape & 24 & 4 & 0.323 & 0.987 & 0.646 \\
summary & 27 & 2 & 0.321 & 0.986 & 0.321 \\
\textbf{behavioral} & \textbf{51} & \textbf{3} & 0.312 & 0.988 & \textbf{0.614} \\
combined & 58 & 3 & 0.279 & 0.983 & 0.623 \\
\bottomrule
\end{tabular}
\end{table}

\textbf{Critical Finding:} Scale-only features achieve highest silhouette (0.521) but \emph{zero} archetype recovery (ARI $-0.004$). This demonstrates internal quality metrics alone cannot validate behavioral segmentation. High silhouette on magnitude simply confirms consumers with different totals are well-separated---but this is the wrong research question.

Behavioral and combined feature sets achieve similar recovery (0.614, 0.623), with behavioral slightly simpler (51 vs 58 features). Shape-only achieves 0.646, suggesting raw hourly bins carry strong signal, but behavioral descriptors add interpretability.

\subsection{Seed Robustness Study}

Twenty independent 30-day datasets (seeds 1-19, 42), five feature sets each:

\begin{table}[h]
\centering
\small
\caption{Robustness Across 20 Datasets (Mean $\pm$ SD)}
\label{tab:robustness}
\begin{tabular}{@{}lccc@{}}
\toprule
Feature Set & Archetype ARI & Silhouette & K Modal (Range) \\
\midrule
\textbf{behavioral} & \textbf{0.641 $\pm$ 0.115} & 0.311 & 3 (3-4) \\
summary & 0.610 $\pm$ 0.240 & 0.319 & 3 (2-4) \\
combined & 0.601 $\pm$ 0.032 & 0.279 & 3 (3-3) \\
shape & 0.589 $\pm$ 0.073 & 0.317 & 3 (3-5) \\
scale & 0.013 $\pm$ 0.036 & \textbf{0.521} & 2 (2-6) \\
\bottomrule
\end{tabular}
\end{table}

\textbf{Statistical Tests (Paired Permutation):}
\begin{itemize}
    \item Behavioral vs scale: $p = 1.9 \times 10^{-6}$ (highly significant)
    \item Behavioral vs shape: $p = 0.123$ (not significant)
    \item Behavioral vs summary: Close, indistinguishable
\end{itemize}

\textbf{Interpretation:} Top three feature sets (behavioral, summary, shape) statistically indistinguishable in recovery. Behavioral selected as shipped feature set based on best mean ARI (0.641) with reasonable interpretability story. Honest reporting: behavioral is not uniquely optimal, but provides best mean performance with scale-invariant design.

Scale features consistently fail ($\text{ARI} \approx 0.01$) despite highest silhouette (0.521), confirming magnitude-timing separation is critical for behavioral segmentation.

\section{Discussion}
\label{sec:discussion}

\subsection{Principal Findings}

This study demonstrates that shape-normalized behavioral features enable household energy clustering by temporal usage patterns independent of consumption magnitude. Four principal findings emerge:

\textbf{1. Magnitude-Timing Separation is Achievable and Critical.} The ablation study (Table~\ref{tab:ablation}) provides conclusive evidence: magnitude-only features achieve near-zero behavioral recovery (ARI $-0.004$) despite highest internal quality (silhouette 0.521). This confirms that raw consumption scale dominates distance calculations when features are not explicitly normalized. The 51-feature behavioral representation successfully isolates timing from magnitude, achieving ARI 0.813 on the flagship 365-day experiment.

\textbf{2. Evidence-Based K Selection Aligns with Ground Truth.} Algorithm~\ref{alg:kselection} selects K=4 using only internal metrics (silhouette, Calinski-Harabasz, Davies-Bouldin, stability), without seeing hidden archetypes. Independent validation confirms K=4 maximizes ARI (0.813) and NMI (0.828) against ground truth. This alignment suggests the composite multi-criterion approach with stability constraints effectively identifies meaningful structure.

\textbf{3. Behavioral Groupings Are Robust Across Time and Initializations.} Multi-seed stability (ARI 0.995) and temporal stability (mean ARI 0.882 across quarters) demonstrate clusters are reproducible properties of consumers, not random artifacts or seasonal effects. Despite 47\% seasonal magnitude variation (Jun 39.4 kWh/day vs Dec 25.3), temporal patterns remain coherent, validating the magnitude-timing separation.

\textbf{4. Internal Metrics Alone Cannot Validate Behavioral Segmentation.} The scale-only arm achieves highest silhouette (0.521) yet zero behavioral recovery, demonstrating that internal quality metrics measure cluster compactness but not semantic meaning. Ground truth validation (even synthetic) or domain expert review remains essential for confirming behavioral interpretability.

\subsection{Comparison with Prior Work}

Our framework extends prior household energy clustering studies \cite{kwac2014household,rhodes2014clustering,jin2016loadshape,lazar2021investigating} through systematic integration of: (1) comprehensive 51-feature behavioral representation beyond simple normalization \cite{valdes2021smart} or raw sequences \cite{kwac2014household}; (2) evidence-based K selection combining multiple metrics with stability rather than heuristic choice; (3) controlled validation quantifying recovery against hidden ground truth; (4) comprehensive stability analysis across both random seeds and temporal windows; (5) rigorous ablation demonstrating magnitude-timing separation; (6) complete reproducibility infrastructure.

Afzalan et al.~\cite{afzalan2020twostage} proposed two-stage clustering separating shape from magnitude using DTW, but did not systematically assess temporal stability or provide controlled validation. Vald\'es et al.~\cite{valdes2021smart} normalized load shapes for demand response but focused on k-medoids without PCA or multi-criterion K selection. Lazar et al.~\cite{lazar2021investigating} identified temperature as dominant variability driver but did not validate clustering stability across time windows.

Our contribution is methodological integration and reproducibility rather than algorithmic novelty in PCA \cite{jolliffe2002pca} or K-means \cite{lloyd1982least,arthur2007kmeans}.

\subsection{Interpretation of Recovered Clusters}

The four clusters exhibit semantically coherent behavioral patterns:

\textbf{Cluster 0 (Midday-Peaking Weekday-Heavy, 19.5\%):} Afternoon peak (13:00), elevated afternoon share (0.357 vs 0.290), depressed weekend ratio (0.747). Consistent with home-office workers, daytime HVAC use, or retirees with daytime routines. Weekday focus suggests occupancy-driven consumption.

\textbf{Cluster 1 (Flat All-Day, 26.0\%):} Lowest peak-to-average (4.92), lowest CV (0.302), near-uniform 24-hour distribution. Resembles commercial-like or always-on baseline consumption. May represent households with minimal occupancy variation or continuous appliance operation.

\textbf{Cluster 2 (Evening-Peaking, 23.5\%):} Highest evening share (0.380), highest peak-to-average (11.32), concentrated 18:00-21:00 demand. Archetypal commuter pattern: minimal daytime use, sharp evening peak for cooking, entertainment, HVAC. Strongest candidate for time-of-use rate targeting.

\textbf{Cluster 3 (Evening-Peaking Weekend-Heavy, 31.0\%):} Evening peak with elevated weekend ratio (1.313). Lifestyle-driven weekly variation: higher weekend consumption suggests recreational home use, social activities, or delayed tasks. Distinct from Cluster 2 by weekend behavior despite similar evening timing.

These interpretations are grounded in SHAP analysis (Table~\ref{tab:shap_features}), which confirms each cluster distinguished by semantically appropriate features (e.g., weekend cluster driven by \texttt{weekend\_ratio} with SHAP importance 0.250).

\subsection{Practical Implications for Demand Response}

Behavioral segmentation enables targeted demand response strategies:

\textbf{Time-of-Use Pricing:} Clusters 2 and 3 (combined 54.5\%) exhibit strong evening peaking, making them prime candidates for evening off-peak incentives. Cluster 0 (19.5\%) midday peakers could benefit from midday solar generation incentives.

\textbf{Load Shifting Programs:} Cluster 2's concentrated evening demand (peak-to-average 11.32) suggests high flexibility potential for pre-cooling, appliance scheduling, or EV charging shifts. Cluster 1's flat profile (peak-to-average 4.92) indicates limited shifting opportunity without baseline reduction.

\textbf{Weekend Programs:} Clusters 0 and 3 exhibit opposite weekend behavior (ratios 0.747 vs 1.313), enabling weekend-specific program targeting. Weekend-heavy consumers (Cluster 3, 31\%) may respond to weekend time-of-use rates differently than weekday-heavy (Cluster 0, 19.5\%).

\textbf{Program Enrollment:} Rather than enrolling all consumers in uniform programs, utilities can design cluster-specific strategies: evening peakers receive evening shift incentives, flat consumers receive baseline reduction programs, weekend-heavy consumers receive weekend-specific rates.

\subsection{Methodological Contributions}

\subsubsection{Shape Normalization}

The explicit separation of magnitude from timing through Equation~\ref{eq:shape_normalization} (section text) addresses a fundamental challenge in energy clustering. Prior work often used raw kWh features \cite{kwac2014household,rhodes2014clustering}, allowing scale to dominate. Simple normalization (z-score) preserves relative magnitude differences. Our approach---dividing by daily total---removes magnitude entirely, enabling pure behavioral clustering. The ablation study conclusively validates this design choice.

\subsubsection{Evidence-Based K Selection}

Algorithm~\ref{alg:kselection} synthesizes best practices from cluster validation literature \cite{hennig2020comparing,gates2019adjusted}: multiple complementary metrics (silhouette, CH, DB), balance constraints (no outlier isolation), stability validation (multi-seed ARI), and parsimony preference (simplest within tolerance). The alignment between selected K=4 and ground truth recovery peak demonstrates this composite approach outperforms single-metric heuristics.

\subsubsection{Controlled Validation}

Synthetic data with hidden archetypes enables quantitative recovery assessment impossible with real unlabeled data. Critical design: magnitude drawn independently of archetype, enabling validation of magnitude-timing separation. This controlled approach complements (not replaces) real-world validation, providing mechanistic understanding of method performance.

\subsubsection{Temporal Stability}

Most energy clustering studies report single-window results without temporal validation \cite{kwac2014household,rhodes2014clustering}. Our quarterly re-analysis with full pipeline re-execution (features, PCA, K-selection) per segment provides stronger evidence: mean ARI 0.882 demonstrates consumer groups are stable longitudinal properties, not seasonal artifacts. This finding has practical importance---consumer behavioral types persist, enabling long-term program targeting.

\subsection{Limitations and Generalization}

See Section~\ref{sec:limitations} for comprehensive discussion. Key constraints:
\begin{itemize}
    \item Primary validation uses synthetic data with simplified archetypes
    \item Real-world pathway implemented but not executed in this study
    \item Moderate scale (200 consumers) vs utility-scale thousands
    \item K-means assumptions (spherical clusters, Euclidean distance)
    \item PCA linearity assumptions
\end{itemize}

Generalization to real smart meter data requires: (1) executing real-world pathway with UCI or utility datasets; (2) domain expert validation of cluster interpretations; (3) assessment at utility scale; (4) comparison with alternative clustering methods (GMM, DBSCAN, hierarchical).

\subsection{Reproducibility as Research Contribution}

Beyond methodological findings, this work demonstrates reproducibility standards for computational energy analytics: deterministic configuration with cryptographic hashing (\texttt{99c7a6631340d301}), pinned dependencies, versioned artifacts, documented execution protocols, public code/data availability, and automated integrity validation. The 0.995 multi-seed stability and 0.882 temporal stability would be undetectable without explicit measurement. Future energy clustering studies should adopt similar reproducibility practices to enable independent verification and meta-analysis.

\section{Limitations}
\label{sec:limitations}

This section explicitly acknowledges methodological constraints, dataset scope, and generalization boundaries. Scientific integrity requires honest limitation reporting rather than concealment.

\subsection{Dataset Limitations}

\subsubsection{Synthetic Primary Dataset}

The flagship validation uses synthetic data generated with four simplified behavioral archetypes (daytime, evening, flat, weekend). While controlled validation enables quantitative recovery assessment (ARI 0.813), it does not prove generalization to real household electricity consumption, which exhibits:
\begin{itemize}
    \item Greater archetype diversity (10-20+ distinct patterns)
    \item Continuous behavior gradients rather than discrete types
    \item Appliance-specific patterns (HVAC, EV charging, pool pumps)
    \item Stochastic occupancy variations
    \item Measurement errors and data quality issues
\end{itemize}

\textbf{Implication:} Real-world validation with utility smart meter data is essential before claiming practical deployment effectiveness.

\subsubsection{Real-World Pathway Not Executed}

A documented pathway for real-world data exists (\texttt{src/dataset\_adapter.py}, \texttt{src/realworld\_ingest.py}, \texttt{src/realworld\_validate.py}) with UCI Individual Household dataset adapter \cite{uci_household}. However, this pathway is \textbf{not executed} in the present study. Claims of real-world validation would be fabricated. The pathway is implementation-ready but requires:
\begin{itemize}
    \item Dataset acquisition and ethics approval
    \item Privacy compliance (anonymization, aggregation)
    \item Domain expert interpretation of clusters
    \item Comparison with utility segmentation practices
\end{itemize}

\textbf{Recommendation:} Execute real-world pathway on UCI dataset or utility partnership data as immediate future work before journal submission.

\subsubsection{Moderate Scale}

The flagship configuration uses 200 consumers over 365 days (1.75 million records). Utility-scale deployments involve:
\begin{itemize}
    \item Thousands to millions of consumers
    \item Multi-year longitudinal data
    \item Heterogeneous building types, climates, demographics
\end{itemize}

\textbf{Scalability Unknown:} Computational complexity (PCA: $O(Np^2)$, K-means: $O(NKmI)$) is feasible for moderate N, but has not been stress-tested at utility scale. Feature engineering on millions of consumers may require distributed computation.

\subsection{Methodological Limitations}

\subsubsection{K-Means Assumptions}

K-means assumes:
\begin{itemize}
    \item Spherical clusters (equal variance in all directions)
    \item Euclidean distance appropriate for similarity
    \item Convex cluster boundaries
    \item Hard assignment (each consumer exactly one cluster)
\end{itemize}

\textbf{Reality:} Behavioral patterns may exhibit:
\begin{itemize}
    \item Elongated or non-spherical distributions
    \item Overlapping gradients (soft memberships)
    \item Non-convex structures
\end{itemize}

\textbf{Alternatives Not Evaluated:} Gaussian Mixture Models (soft assignment, elliptical clusters), DBSCAN (density-based, arbitrary shapes), hierarchical clustering (nested structures) may better capture complex behavioral distributions. Benchmark comparison is future work.

\subsubsection{PCA Linearity}

PCA assumes linear relationships between features. Behavioral patterns may involve:
\begin{itemize}
    \item Nonlinear interactions (e.g., quadratic temperature effects)
    \item Multiplicative relationships
    \item Threshold effects
\end{itemize}

\textbf{Alternative:} Kernel PCA, autoencoders, or nonlinear manifold learning (t-SNE, UMAP) may capture richer structure but sacrifice interpretability and increase computational cost.

\subsubsection{Feature Engineering as Inductive Bias}

The 51 behavioral features encode domain knowledge and design choices:
\begin{itemize}
    \item Period definitions (4 blocks of 6 hours)
    \item Harmonic counts (3 components)
    \item Haar wavelet levels (3)
    \item Weekend definition (Saturday-Sunday)
\end{itemize}

\textbf{Limitation:} Different feature sets may yield different clusterings. The ablation study shows behavioral features outperform alternatives, but optimal feature engineering remains an open question. Learned representations (deep learning) may discover features humans miss.

\subsubsection{Single Flagship Run Without Uncertainty}

The 365-day flagship experiment is executed once (seed 42). While multi-seed stability at K=4 is high (ARI 0.995), we do not report:
\begin{itemize}
    \item Confidence intervals on ARI/NMI
    \item Bootstrap uncertainty on cluster profiles
    \item Sensitivity to small perturbations
\end{itemize}

\textbf{Robustness study} (20 datasets, 30-day) provides mean/variance estimates but on shorter window. Flagship uncertainty quantification is future work (bootstrap resampling recommended).

\subsection{Validation Limitations}

\subsubsection{Internal Metrics Are Heuristics}

Silhouette, Calinski-Harabasz, and Davies-Bouldin measure cluster compactness and separation but do not confirm:
\begin{itemize}
    \item Semantic meaningfulness
    \item Actionability for applications
    \item Stability under real-world conditions
\end{itemize}

The scale-only ablation arm (silhouette 0.521, ARI $-0.004$) demonstrates internal metrics can mislead. \textbf{Domain validation essential.}

\subsubsection{Ground Truth Limited to Synthetic Data}

ARI=0.813 recovery validates the method \emph{on synthetic data with known archetypes}. Real households have no ground truth labels, so external validation metrics (ARI, NMI) are unavailable. Real-world validation must rely on:
\begin{itemize}
    \item Domain expert assessment
    \item Comparison with utility segmentation
    \item Intervention studies (demand response trials)
\end{itemize}

\subsubsection{Explainability is Post-Hoc}

SHAP analysis explains a \emph{surrogate} random forest (CV accuracy 0.985), not the clustering directly. While high surrogate accuracy suggests features strongly distinguish clusters, the explanation is indirect. Clustering has no native feature importance, and SHAP values are model-dependent.

\subsection{Temporal Scope}

\subsubsection{Horizon Dependence}

Shorter windows under-recover archetypes: 30-day reference selects K=3 (ARI 0.61) vs 365-day K=4 (ARI 0.81). \textbf{Implication:} Behavioral segmentation requires sufficient observation time to capture weekly and seasonal variations. Utilities implementing clustering should use $\geq$180 days (documented threshold for longitudinal analysis).

\subsubsection{Single-Year Data}

Flagship uses 2024 only. Multi-year patterns (e.g., household lifecycle changes, appliance upgrades) are not assessed. Longitudinal cluster membership stability over years is unknown.

\subsection{Generalization Boundaries}

\textbf{What This Study Demonstrates:}
\begin{itemize}
    \item Shape normalization successfully separates magnitude from timing
    \item Evidence-based K selection aligns with ground truth on synthetic data
    \item Behavioral features achieve high recovery (ARI 0.813)
    \item Clusters are stable across seeds (ARI 0.995) and quarters (ARI 0.882)
    \item Magnitude-only features fail behavioral recovery (ARI $-0.004$)
\end{itemize}

\textbf{What This Study Does NOT Demonstrate:}
\begin{itemize}
    \item Real-world clustering effectiveness (pathway not executed)
    \item Generalization to diverse climates, building types, demographics
    \item Utility-scale computational feasibility (tested at N=200)
    \item Superiority over alternative methods (no benchmark comparison)
    \item Demand response intervention effectiveness (observational study only)
    \item Cost-effectiveness or return on investment
\end{itemize}

\subsection{Ethical and Privacy Considerations}

While this study uses synthetic data, real-world deployment raises:
\begin{itemize}
    \item \textbf{Privacy:} Smart meter data reveals occupancy, appliance use, lifestyle patterns. Clustering amplifies re-identification risk through behavioral fingerprinting. Utilities must implement privacy-preserving aggregation, differential privacy, or secure multi-party computation.
    \item \textbf{Equity:} Behavioral segmentation may create disparate impacts if clusters correlate with income, demographics, or vulnerability. Time-of-use pricing favoring flexible consumers may burden shift workers, caregivers, or those with inflexible schedules.
    \item \textbf{Consent:} Consumers may not understand clustering implications. Transparent communication and opt-out mechanisms are essential.
\end{itemize}

\subsection{Addressing Limitations in Future Work}

\textbf{High Priority:}
\begin{enumerate}
    \item Execute real-world pathway on UCI dataset or utility partnership
    \item Add bootstrap confidence intervals for flagship metrics
    \item Compare with GMM, DBSCAN, hierarchical clustering
\end{enumerate}

\textbf{Medium Priority:}
\begin{enumerate}
    \item Scale to utility-scale thousands of consumers
    \item Multi-year longitudinal validation
    \item Intervention study (demand response trial)
\end{enumerate}

\textbf{Low Priority:}
\begin{enumerate}
    \item Kernel PCA or autoencoder alternatives
    \item Hyperparameter sensitivity analysis
    \item Feature selection optimization
\end{enumerate}

\subsection{Conclusion on Limitations}

These limitations do not invalidate the methodological contributions but establish boundaries on claims. This is a \textbf{methodology paper} demonstrating reproducible behavioral segmentation on controlled synthetic data. Real-world effectiveness requires additional validation before deployment recommendations.

\section{Conclusion}
\label{sec:conclusion}

This paper presents a reproducible framework for shape-first behavioral segmentation of household energy consumption, addressing the fundamental challenge of separating temporal usage patterns from consumption magnitude. Through controlled synthetic validation, comprehensive stability analysis, rigorous ablation studies, and complete reproducibility infrastructure, we demonstrate that behavioral clustering by \emph{when} energy is consumed (not \emph{how much}) is achievable and actionable for demand response applications.

\subsection{Key Contributions}

\textbf{1. Methodological Integration:} We synthesize shape normalization, 51-feature behavioral representation, PCA dimensionality reduction, evidence-based K selection, and multi-axis stability validation into a cohesive, reproducible pipeline. While individual techniques are established (PCA \cite{jolliffe2002pca}, K-means \cite{lloyd1982least,arthur2007kmeans}), their rigorous integration with controlled validation and comprehensive stability analysis advances methodological standards for behavioral energy analytics.

\textbf{2. Empirical Validation of Magnitude-Timing Separation:} The ablation study provides conclusive evidence: magnitude-only features achieve highest internal quality (silhouette 0.521) yet zero behavioral recovery (ARI $-0.004$), while shape-normalized behavioral features achieve ARI 0.813. This demonstrates that internal metrics alone cannot validate behavioral segmentation---external validation (even synthetic ground truth) is essential.

\textbf{3. Evidence-Based K Selection:} Algorithm~\ref{alg:kselection} integrates multiple complementary metrics (silhouette, Calinski-Harabasz, Davies-Bouldin), balance constraints, stability validation, and parsimony preference. The selected K=4 independently aligns with ground truth recovery peak (ARI 0.813), validating the composite approach over single-metric heuristics.

\textbf{4. Comprehensive Stability Analysis:} Multi-seed stability (ARI 0.995) and temporal stability across quarters (mean ARI 0.882) demonstrate clusters are reproducible and persist across seasons despite 47\% magnitude variation. This longitudinal validation, rarely reported in energy clustering literature, confirms behavioral types are stable consumer properties.

\textbf{5. Reproducibility Infrastructure:} Deterministic configuration with cryptographic hashing, pinned dependencies, versioned artifacts, automated integrity checks, and public code/data availability establish reproducibility standards. The framework enables independent verification, sensitivity analysis, and extension by other researchers.

\subsection{Practical Implications}

Behavioral segmentation enables targeted demand response strategies: evening-peaking clusters (54.5\% of flagship consumers) can receive time-of-use pricing incentives; flat-profile consumers (26\%) require baseline reduction programs; weekend-heavy consumers (31\%) may respond to weekend-specific rates. Rather than uniform enrollment, utilities can design cluster-specific interventions, potentially improving program effectiveness and consumer satisfaction.

However, practical deployment requires real-world validation (Section~\ref{sec:limitations}). The documented UCI household adapter pathway provides implementation-ready infrastructure, but execution with utility partnership data is essential before operational recommendations.

\subsection{Scientific Positioning}

This work is a \textbf{methodology paper} contributing reproducible research practices, not algorithmic novelty. PCA and K-means are well-established; our contribution is rigorous application with controlled validation, comprehensive stability analysis, and reproducibility infrastructure. We emphasize what has been demonstrated (synthetic data recovery, stability, magnitude-timing separation) and explicitly acknowledge what has not (real-world validation, utility-scale feasibility, benchmark superiority).

\subsection{Limitations and Boundaries}

Key constraints include: synthetic primary dataset with simplified archetypes, real-world pathway implemented but not executed, moderate scale (200 consumers), K-means spherical cluster assumptions, PCA linearity assumptions, single flagship run without bootstrap confidence intervals, and absence of benchmark comparison with alternative clustering methods (GMM, DBSCAN, hierarchical). These limitations establish boundaries on claims but do not invalidate methodological contributions.

\subsection{Future Research Directions}

\textbf{Immediate Priorities:}
\begin{enumerate}
    \item Execute real-world UCI pathway with domain expert validation
    \item Add bootstrap confidence intervals for flagship uncertainty quantification
    \item Benchmark against Gaussian Mixture Models, DBSCAN, hierarchical clustering
    \item Scale to utility-level thousands of consumers
\end{enumerate}

\textbf{Methodological Extensions:}
\begin{enumerate}
    \item Nonlinear dimensionality reduction (kernel PCA, autoencoders)
    \item Soft clustering with probabilistic assignments (GMM, fuzzy c-means)
    \item Hierarchical behavioral taxonomies
    \item Multi-resolution temporal features (hourly, daily, weekly, seasonal)
    \item Transfer learning across utilities and climates
\end{enumerate}

\textbf{Application Studies:}
\begin{enumerate}
    \item Demand response intervention trials with randomized cluster assignment
    \item Cost-effectiveness analysis of cluster-specific vs. uniform programs
    \item Longitudinal cluster membership stability over multi-year periods
    \item Privacy-preserving clustering with differential privacy
    \item Equity analysis of behavioral segmentation impacts
\end{enumerate}

\subsection{Broader Impact}

Beyond household energy, shape-normalized behavioral clustering generalizes to:
\begin{itemize}
    \item Commercial building energy management
    \item Water consumption pattern analysis
    \item Transportation mobility segmentation
    \item Healthcare activity pattern recognition
    \item Network traffic classification
\end{itemize}

The methodological framework---separate magnitude from behavior, apply evidence-based selection, validate through multiple axes, ensure reproducibility---applies wherever temporal patterns must be distinguished from scale.

\subsection{Reproducibility and Open Science}

All materials are publicly available:
\begin{itemize}
    \item Repository: \url{https://github.com/shaxntanu/Energy-Consumption-Pattern-Analysis-using-PCA-and-K-Means}
    \item Interactive Explorer: \url{https://energy-consumption-pattern.vercel.app}
    \item Configuration Hash: \texttt{99c7a6631340d301} (flagship 365-day run)
    \item Zenodo DOI: [To be assigned upon publication]
\end{itemize}

We encourage researchers to:
\begin{itemize}
    \item Replicate experiments using documented protocols
    \item Extend framework to real-world datasets
    \item Compare with alternative clustering approaches
    \item Adapt to other temporal behavioral analysis domains
    \item Report reproducibility metrics (configuration hashes, stability, temporal validation)
\end{itemize}

\subsection{Final Remarks}

Household energy consumption clustering has traditionally focused on magnitude-based segmentation, grouping consumers by \emph{how much} they use. This paper demonstrates that behavioral segmentation by \emph{when} they use energy is achievable through shape normalization, provides complementary insights for demand response, and requires fundamentally different validation approaches. Internal quality metrics (silhouette, Calinski-Harabasz, Davies-Bouldin) measure cluster compactness but cannot confirm behavioral meaningfulness. External validation---whether through synthetic ground truth, domain experts, or intervention trials---is essential.

The methodological contributions (shape normalization, evidence-based K selection, comprehensive stability analysis, reproducibility infrastructure) advance standards for computational energy analytics. Real-world validation remains critical future work. We have established methodological rigor on controlled synthetic data; the next step is demonstrating effectiveness on utility smart meter deployments.

Reproducible research requires honest limitation reporting, independent verification, and transparent uncertainty quantification. This paper aims to set standards: what has been demonstrated, what has not, what remains unknown. We hope future energy clustering studies adopt similar practices, enabling cumulative scientific progress through independently verifiable results.


% ============================================================================
% ACKNOWLEDGMENTS
% ============================================================================
\section*{Acknowledgments}
The authors acknowledge the use of the Zephyr Station weather API for seasonal data integration. This research was conducted as part of an academic project on energy consumption pattern analysis.

% ============================================================================
% DATA AVAILABILITY
% ============================================================================
\section*{Data Availability}
The synthetic dataset generator, complete analysis pipeline, experimental configurations, and all results are publicly available at: \url{https://github.com/shaxntanu/Energy-Consumption-Pattern-Analysis-using-PCA-and-K-Means}. Configuration hash for flagship run: \texttt{99c7a6631340d301}. Interactive results explorer: \url{https://energy-consumption-pattern.vercel.app}.

% ============================================================================
% CODE AVAILABILITY
% ============================================================================
\section*{Code Availability}
All source code is released under MIT License. Core modules: data generation (\texttt{src/data\_loader.py}), preprocessing (\texttt{src/preprocessing.py}), feature engineering (\texttt{src/feature\_engineering.py}), PCA (\texttt{src/pca\_analysis.py}), clustering (\texttt{src/clustering.py}), validation (\texttt{src/validation.py}), and stability analysis (\texttt{src/longitudinal\_analysis.py}, \texttt{src/seasonal\_analysis.py}). Complete test suite available in \texttt{tests/}.

% ============================================================================
% BIBLIOGRAPHY
% ============================================================================
\printbibliography

\end{document}
